El desarrollo de la práctica se rigió bajo estrictos principios de aislamiento y ética profesional. La arquitectura del sistema se diseñó mediante un esquema de red local (LAN) aislada, cumpliendo con las normativas institucionales de la universidad. Esta configuración permitió mitigar los riesgos asociados a la exposición de direcciones IP públicas, garantizando que el tráfico malicioso generado se mantuviera en un entorno controlado y seguro, sin vulnerar la integridad de la red externa.\\
Complementariamente, se aplicó una metodología de observación pasiva. A diferencia de los sistemas de prevención tradicionales, el entorno se configuró deliberadamente para no bloquear las intrusiones de forma inmediata. El objetivo de este enfoque fue permitir una interacción extendida de los atacantes con los servicios simulados, facilitando la recolección de evidencia y el estudio detallado de sus técnicas de explotación (como el análisis de payloads y comandos ejecutados) sin comprometer activos ni recursos reales de la organización.